\section{An actual Euler commutator in Meyer's common-space trace}
\label{sec:meyer-transport}

This section supplies a prime-carrying operator calculation on the
one-sided completions used by Meyer.  It neither assumes a zero-location
theorem nor extends the canonical inverse boundedly to the matching
narrow two-sided norms refuted in MA-MO-001.  The prime contribution is
the trace of the indicated operator, obtained from its kernel.  No
operator is defined by prescribing the value of its trace.

We use $x=\log t$, convolution
$C_g u(x)=\int_{\mathbb R}g(y)u(x-y)\,dy$, and
$T_\ell u(x)=u(x+\ell)$.  Let $S$ be a finite set of rational primes,
let $\mathcal N_S$ be the positive integers whose prime factors are
outside $S$, and set
\[
 \mathcal H_a=L^2(\mathbb R,e^{2ax}\,dx),\qquad a>1/2,
 \qquad
 Z_S=\sum_{n\in\mathcal N_S}n^{-1/2}T_{\log n},\quad
 J_S=\sum_{n\in\mathcal N_S}\mu(n)n^{-1/2}T_{\log n}.
 \tag{MA-MEY.1}
\]
Here $J_S$ denotes the M\"obius inverse.  Meyer's letter $J$ denotes
an inversion operator; we denote that different operator by $J_0$
below.  The notation is deliberately kept separate.

\paragraph{MA-MEY-001 (a smoothed Euler commutator and its actual trace).}
Both series in (MA-MEY.1) converge absolutely in operator norm,
and $Z_SJ_S=J_SZ_S=1$.  Let $g\in C_c^\infty(\mathbb R)$ and let
$\rho\in C^\infty(\mathbb R)$ have compactly supported derivative,
with limiting constants $\rho_-$ at $-\infty$ and $\rho_+$ at
$+\infty$.  Then
\[
 K_{S,\rho}(g)=C_gZ_S[J_S,M_\rho]
             =C_g(M_\rho-Z_SM_\rho J_S)
 \tag{MA-MEY.2}
\]
is trace class on $\mathcal H_a$, and
\[
 \boxed{\operatorname{Tr}_{\mathcal H_a}K_{S,\rho}(g)
 = (\rho_- -\rho_+)
   \sum_{\substack{p\notin S\\ k\ge1}}
     (\log p)p^{-k/2}g(k\log p).}
 \tag{MA-MEY.3}
\]
The prime sum is finite.  The trace-norm bound is uniform for $g$
in a bounded subset of $C_c^\infty$ with fixed compact support.
There is also a trace-class formula before multiplication by $Z_S$:
\[
 \operatorname{Tr}_{\mathcal H_a}C_g[J_S,M_\rho]
 = (\rho_+-\rho_-)
   \sum_{n\in\mathcal N_S}\mu(n)(\log n)n^{-1/2}g(\log n).
 \tag{MA-MEY.4}
\]

\begin{proof}
Since $\|T_{\log n}\|_{\mathcal H_a}=n^{-a}$, both series are
absolutely convergent for $a>1/2$.  Their product is an absolutely
convergent double series.  Its coefficient at $T_{\log r}$ is
$r^{-1/2}\sum_{d\mid r}\mu(d)$ for $r\in\mathcal N_S$.
This is one at $r=1$ and zero otherwise, proving the inverse
identities.  Compact smooth convolution is bounded: conjugating by
the unitary $W_a u(x)=e^{ax}u(x)$ gives
$W_aC_gW_a^{-1}=C_{g_a}$, where $g_a(y)=e^{ay}g(y)$, and
Young's inequality applies.

We give a trace-norm estimate, not merely a formal diagonal
calculation.  If $q\in C_c^\infty(\mathbb R)$ and
$b\in C_c^\infty(\mathbb R)$ has support in an interval of length
$L$, then
\[
 \|C_qM_b\|_1\le
 c_q(1+L)\max_{0\le j\le2}\|b^{(j)}\|_\infty.
 \tag{MA-MEY.5}
\]
To prove this, choose a fixed smooth partition of unity
$\sum_{r\in\mathbb Z}\eta(x-r)=1$ with compactly supported
$\eta$.  Only $O(1+L)$ of the functions
$b_r(x)=\eta(x-r)b(x)$ are nonzero.  Simultaneously translating
both variables reduces the kernel
$q(x-z)b_r(z)$ to a smooth kernel in a fixed compact rectangle,
whose size depends only on $\operatorname{supp}q$ and $\eta$.
Choose a larger fixed square with the rectangle in its interior
and expand the kernel in the orthonormal Fourier basis on that
square, extended by zero to $\mathbb R$.  Integrating twice in
each variable bounds its coefficients by
\[
 c_q\max_{0\le j\le2}\|b^{(j)}\|_\infty
 (1+m^2)^{-1}(1+n^2)^{-1}.
\]
Only derivatives of $q$ through order four and of $b_r$ through
order two occur.  The absolute sum of these coefficients bounds
the trace norm of the corresponding rank-one series.  Summing
over $r$ proves (MA-MEY.5).  The same absolute expansion, whose
diagonal terms are integrable by Cauchy--Schwarz, proves that the
trace of each of these compact smooth kernels is its diagonal
integral.

For $\ell\ge0$, put
$D_\ell(x)=\rho(x+\ell)-\rho(x)$.  If
$\operatorname{supp}\rho'\subseteq[A,B]$, then
$\operatorname{supp}D_\ell\subseteq[A-\ell,B]$;
all its derivatives through order two are bounded uniformly in
$\ell$.  The identity
\[
 [T_\ell,M_\rho]=M_{D_\ell}T_\ell
\]
and (MA-MEY.5), after conjugation by $W_a$, imply
\[
 \|C_g[T_\ell,M_\rho]\|_{1,\mathcal H_a}
       \le c_{g,a,\rho}e^{-a\ell}(1+\ell).
 \tag{MA-MEY.6}
\]
Here the final translation is unitary on ordinary $L^2$, and
$W_aT_\ell W_a^{-1}=e^{-a\ell}T_\ell$.
Thus the series for $C_g[J_S,M_\rho]$ converges absolutely in
trace norm.  More precisely, the double series
\[
 K_{S,\rho}(g)
  =\sum_{m,n\in\mathcal N_S}
       \mu(n)(mn)^{-1/2}
       C_gT_{\log m}[T_{\log n},M_\rho]
 \tag{MA-MEY.7}
\]
does so, since the trace norm of the term without its scalar
coefficient is at most
$c_{g,a,\rho}(mn)^{-a}(1+\log n)$.  This follows from
(MA-MEY.5) also when $D_{\log n}$ is translated by $\log m$.
The convergent majorant is
\[
 c_{g,a,\rho}
 \sum_{m\ge1}m^{-a-1/2}
 \sum_{n\ge1}n^{-a-1/2}(1+\log n).
\]
The constants use only a fixed finite list of seminorms of $g$
on its support.  This proves the stated uniformity.

Write $v=\log m$ and $w=\log n$.  The kernel with respect to
ordinary $dz$ of the term $C_gT_v[T_w,M_\rho]$ is
\[
     g(x+v+w-z)\{\rho(z)-\rho(z-w)\}.
 \tag{MA-MEY.8}
\]
Its conjugate by $W_a$ has the additional factor $e^{a(x-z)}$,
which is one on the diagonal.  Equivalently, its diagonal
factor is $e^{-a(v+w)}g_a(v+w)=g(v+w)$.  The compact-kernel
trace calculation just proved therefore gives
\[
 \operatorname{Tr} C_gT_v[T_w,M_\rho]
   =g(v+w)\int_{\mathbb R}\{\rho(x)-\rho(x-w)\}\,dx
   =g(v+w)w(\rho_+-\rho_-).
 \tag{MA-MEY.9}
\]
For the last equality, write the difference as
$\int_{x-w}^x\rho'(u)\,du$ and use Fubini; the integrand is
absolutely integrable since $\rho'$ has compact support.

Taking traces in the absolutely trace-norm convergent series
(MA-MEY.7), and regrouping by $r=mn$, gives
\[
 \operatorname{Tr}K_{S,\rho}(g)
 = (\rho_+-\rho_-)
   \sum_{r\in\mathcal N_S}r^{-1/2}g(\log r)
              \sum_{n\mid r}\mu(n)\log n.
\]
The inner sum is $-\Lambda(r)$.  Indeed, for a prime power the
only squarefree divisors are $1$ and $p$, giving $-\log p$;
if $r$ has at least two distinct prime factors, the coefficient
of each $\log p$ cancels in the product expansion of
$\prod_{q\mid r}(1-1)$.  It is zero also at $r=1$.
This proves (MA-MEY.3).  Taking only $m=1$ in the kernel
calculation proves (MA-MEY.4).  Compactness of
$\operatorname{supp}g$ makes both final arithmetic sums finite.
\end{proof}

The multiplication by $Z_S$ is essential to the interpretation:
(MA-MEY.4) differentiates the reciprocal Euler multiplier,
whereas (MA-MEY.3) produces its logarithmic derivative.  With
the Mellin convention of MA-TR-101, their multipliers are
$J_S(z)=1/\zeta_{\overline S}(z+1/2)$ and
$J_S'(z)/J_S(z)=-\zeta_{\overline S}'/
\zeta_{\overline S}(z+1/2)$ in $\Re z>1/2$.
The signs in (MA-MEY.3)--(MA-MEY.4) are fixed by the displayed
operator order and the limits of $\rho$.

\paragraph{MA-MEY-002 (a boundary contribution invisible to compact cutoffs).}
If $\rho$ has compact support, then
$\operatorname{Tr}K_{S,\rho}(g)=0$.  If
$\phi\in C^\infty(\mathbb R)$ satisfies $\phi=1$ near
$-\infty$ and $\phi=0$ near $+\infty$, then
\[
 \operatorname{Tr}K_{S,\phi}(g)=P_{\overline S}^+(g),\qquad
 P_{\overline S}^+(g)=
 \sum_{\substack{p\notin S\\k\ge1}}
       (\log p)p^{-k/2}g(k\log p).
 \tag{MA-MEY.10}
\]
There are uniformly bounded $\rho_R\in C_c^\infty$ with
$\rho_R\to\phi$ pointwise, and for every such family
$K_{S,\rho_R}(g)\to K_{S,\phi}(g)$ strongly on $\mathcal H_a$.
Whenever $P_{\overline S}^+(g)\ne0$, this convergence is not in
trace norm: for every $R$,
\[
 \|K_{S,\rho_R}(g)-K_{S,\phi}(g)\|_1
       \ge |P_{\overline S}^+(g)|.
 \tag{MA-MEY.11}
\]

\begin{proof}
The first two claims follow from (MA-MEY.3), since the two
limits of a compact cutoff are zero.  Multiplying $\phi$ by
a smooth function that is zero on $(-\infty,-R-1]$ and one on
$[-R,\infty)$ supplies the stated family.  Uniform boundedness
and dominated convergence give $M_{\rho_R}\to M_\phi$ strongly.
All other operators in (MA-MEY.2) are fixed bounded operators,
so the same is true of $K$.  Each difference is trace class;
the trace inequality $|\operatorname{Tr}T|\le\|T\|_1$ gives
(MA-MEY.11).  For a concrete nonzero case, take $S=\varnothing$
and $g\ge0$ supported sufficiently close to $\log2$ that it
contains no other positive logarithm of a prime power, with
$g(\log2)=1$.  Then
$P_{\overline S}^+(g)=\log2/\sqrt2$.
\end{proof}

This is a precise obstruction to replacing Meyer's one-sided
cutoff by compact cutoffs in this particular trace-class
commutator.  It does not exclude other regularizations.
Nor is there a contradiction with cyclicity of the trace:
$C_gM_\phi$ and $C_gZ_SM_\phi J_S$ have not separately been
shown trace class.  Their difference has the trace just computed.
For compact $\rho$, both separated terms are trace class and
cyclicity indeed gives zero.

\paragraph{MA-MEY-003 (the fixed-$S$ defect in Meyer's actual common space).}
Restrict Meyer's construction over $\mathbb Q$ to the trivial
compact norm-one isotype, with that compact group of volume one.
Use the finite place set $S\cup\{\infty\}$ and a symmetric
cutoff $\phi(x)+\phi(-x)=1$ in logarithmic coordinates, with
the endpoint values in MA-MEY-002.  Let $P_+^S=i_+p_+^S$
and $P_-=i_-p_-$ be Meyer's actual approximate projection and
projection.  On the centered common completion
$\mathcal H_a\oplus\mathcal H_{-a}$, $a>1/2$, their smoothed
difference is trace class, and
\[
 \operatorname{Tr}_{\mathcal H_a\oplus\mathcal H_{-a}}
        C_g(P_- -P_+^S)
 =-B_{\overline S}(g),\qquad
 B_{\overline S}(g)=P_{\overline S}^+(g)
                         +P_{\overline S}^+(g^\vee),
 \quad g^\vee(x)=g(-x).
 \tag{MA-MEY.12}
\]
Here $C_g$ acts diagonally on the common space.  Thus the
omitted-prime correction in Meyer's trace additivity is an
actual trace of the common-space defect, even for a fixed $S$
which has not absorbed every prime meeting the support of $g$.

\begin{proof}
We first record the exact source maps to fix their spaces and
normalizations.  In Meyer's notation let $\mathfrak F$ be the
additive Fourier transform, $J_0 f(t)=t^{-1}f(t^{-1})$ the
inversion, and $\mathsf F=\mathfrak FJ_0$ the equivariant
Fourier transform.  The common smooth space is
$\mathcal S(\mathcal C)_{>}\oplus
 \mathcal S(\mathcal C)_{<}$.
The actual maps are
\begin{align*}
 i_+v&=(\Sigma v,J_0\Sigma\mathfrak F^*v),\\
 p_+^S(v_0,v_1)
   &=M_{1-\phi}\mathcal L_S^{-1}v_0
       +\mathfrak F M_{1-\phi}\mathcal L_S^{-1}J_0v_1,\\
 i_-v&=(v,v),\qquad
 p_-(v_0,v_1)=M_{1-\phi}v_0+M_\phi v_1.
 \tag{MA-MEY.13}
\end{align*}
Here $\Sigma$ is arithmetic summation and
$\mathcal L_S=\prod_{p\notin S}(1-\lambda_p^{-1})^{-1}$
is its realization on the $S$-local space $\mathcal T$.
These statements are Meyer (2005), Lemmas~5.3--5.4,
pp.~39--40, and the displays immediately following Lemma~5.4.
The maps are defined on the local and global smooth spaces
there; the proof of the nuclear defect supplies the indicated
Hilbert extensions, rather than assuming that all projection
entries individually act on that Hilbert direct sum.
For $\mathbb Q$, every set $S\cup\{\infty\}$ here is
``sufficiently large'' in Meyer's structural sense: unique
factorization identifies
$\mathbb Q^\times/\mathbb Z[S^{-1}]^\times$ with the free
valuation group on primes outside $S$. This is the definition
in his \S5.1, before the summation construction; it does not
require enlarging $S$ to fit the test support.

Center functions by $v(t)\mapsto u(x)=e^{x/2}v(e^x)$.
The raw Hilbert weight $\alpha>1$ becomes
$a=\alpha-1/2>1/2$, and its paired weight $1-\alpha$
becomes $-a$.  Inversion becomes the unitary reflection
$Ru(x)=u(-x)$ from $\mathcal H_a$ to $\mathcal H_{-a}$.
The operator $\mathcal L_S$ becomes $Z_S$, its inverse
becomes $J_S$, and raw convolution by
$f(t)=t^{-1/2}g(\log t)$ becomes $C_g$.  These statements
follow directly from $\lambda_t v(u)=v(t^{-1}u)$; in
particular the half-density factor in $\lambda_{1/n}$ is
$n^{-1/2}$, exactly as in (MA-MEY.1).

Meyer's Lemma~5.6, pp.~41--42, proves that the smoothed
difference $C_g(P_- -P_+^S)$ is nuclear from the raw weighted
Hilbert direct sum into the common smooth space.  Composing
with the continuous inclusion into the same Hilbert sum makes
it trace class there. The nuclear trace on the common smooth
space equals this Hilbert trace: one nuclear representation
$\sum_j\lambda_j v_j\otimes\ell_j$ of the Hilbert-to-smooth
map gives both traces as the same absolutely convergent series
$\sum_j\lambda_j\ell_j(v_j)$, with the inclusion understood.
This is precisely the compatible-pairing argument in Meyer's
Corollary \texttt{cor:T\_on\_Twist}; it uses the actual nuclear
lift, not just a range inclusion for a trace-class endomorphism.
This applies in particular on the
invariant compact isotype considered here.  No condition
$q_p>\sup|\operatorname{supp}f|$ is required for that lemma;
the support-adapted enlargement of $S$ is used later to make
the trace vanish.

From the definitions (MA-MEY.13), the two diagonal entries
of the centered difference $P_+^S-P_-$ are $D$ and $RDR$,
where
\[
 D=Z_SM_{1-\phi}J_S-M_{1-\phi}
   =M_\phi-Z_SM_\phi J_S=Z_S[J_S,M_\phi].
 \tag{MA-MEY.14}
\]
For the second diagonal one uses
$RM_{1-\phi}R=M_\phi$ and $Z_SJ_S=1$.
The trace of a trace-class operator on a Hilbert direct sum
is the sum of the traces of its diagonal entries.  Since
$RC_gR=C_{g^\vee}$, (MA-MEY.3) gives
\[
 \operatorname{Tr}C_g(P_+^S-P_-)
 =\operatorname{Tr}_{\mathcal H_a}C_gD
  +\operatorname{Tr}_{\mathcal H_a}C_{g^\vee}D
 =B_{\overline S}(g).
\]
This proves (MA-MEY.12), including its sign.
\end{proof}

There is a printed sign inconsistency that matters for this
nonzero refinement.  In the proof of Meyer's Lemma~5.6,
the display labelled \texttt{eq:global\_perturbation} in the
attached primary \LaTeX\ labels its left side $P_- -P_+^S$
but gives the entries of $P_+^S-P_-$.  The definitions of
$P_+^S$ and $P_-$ just before \S5.3 determine the sign in
(MA-MEY.12).  Meyer's later vanishing argument is unchanged
by this printed sign; no failure of the global trace theorem
is inferred from it.

For completeness, the off-diagonal part of that source proof
also has signs that should be read with the defining Poisson
formula. Here is the short corrected reduction. In raw
coordinates put $q=1-\phi$, $L=\mathcal L_S$ and
$\mathcal C_u=\lambda(u)$ for convolution. Write
$s=\operatorname{Sing}(J_0f)$, so Poisson gives
$LJ_0f=h-\phi s$ with $h\in\mathcal S(\mathcal C_S)_\cup$.
The relevant off-diagonal operator is
\[
 \begin{split}
 O_f&=\mathcal C_f
          (L\mathfrak F M_qL^{-1}-M_\phi J_0)\\
 &=-[\mathcal C_f,M_\phi]J_0
   -J_0\bigl([M_q,\mathcal C_{LJ_0f}]
                         -\mathcal C_sM_q\bigr)L^{-1},\\
 [M_q,\mathcal C_{LJ_0f}]-\mathcal C_sM_q
   &=[M_q,\mathcal C_h]
                  -M_q\mathcal C_{\phi s}-\mathcal C_{qs}M_q.
 \end{split}
\]
These identities follow from $J_0L\mathfrak F=L+\operatorname{Sing}$
and the convolution intertwining relation. The last two kernels
are separately nuclear from the raw $L^2_\alpha$, $\alpha>1$,
to $\mathcal S(\mathcal C_S)_{<}$: $s$ is a linear combination
of $1$ and $t^{-1}$, and their supports respectively force
$1\ll x\ll y$ and $1\ll y\ll x$, up to fixed constants.
After weighting by $x^\beta y^{-\alpha}$, $\beta<0$, both
kernels and every logarithmic derivative decay rapidly.
The other term is the source's nuclear commutator lemma.
Reflection gives the required $>$ target. Thus the corrected
signs preserve the full nuclear conclusion of Lemma~5.6.

\subsection{The completed geometric piece and the remaining transport}

For precision, the trace additivity used in the proof of
Meyer's Theorem~5.8, pp.~43--44, is
\begin{align*}
 \chi_{\pi_+}(f)-\chi_{\pi_-}(f)
 &=\operatorname{Tr}_{\mathcal T}
       \lambda(f)(1-p_+^Si_+)
       -\operatorname{Tr}_{\mathrm{common}}
          \lambda(f)(P_- -P_+^S),\\
 1-p_+^Si_+&=M_\phi-\mathsf F M_\phi\mathsf F^{-1}.
 \tag{MA-MEY.15}
\end{align*}
The quotient $\pi_+$ is the two-dimensional polar quotient;
$\pi_-$ is the negative smooth quotient.  The maps and the
additivity are taken in the actual common space, using Poisson
summation and nuclearity.  The first trace is the actual local
Fourier-cutoff trace.  Meyer proves that its computation on
$\mathcal T$ agrees with its computation on
$L^2(\mathcal C_S)_{1/2}$ in Proposition
\texttt{pro:T\_on\_Twist} and Corollary
\texttt{cor:T\_on\_Twist}; the local value is calculated in
Lemma~3.3 and the following local principal-value calculation.
Combining (MA-MEY.12) and (MA-MEY.15) over $\mathbb Q$ gives
\[
 \chi_{\pi_+}(f)-\chi_{\pi_-}(f)
   =W_S(f)+B_{\overline S}(g),
 \qquad f(t)=t^{-1/2}g(\log t),
 \tag{MA-MEY.16}
\]
where $W_S$ is this independently computed local trace, with
the source's fixed Haar and principal-value conventions.
The discriminant term is zero for $\mathbb Q$.  This proof
uses the source's trace additivity before its final choice of
$S$, and (MA-MEY.3) computes the term that that choice makes
zero.  It does not begin with the numerical explicit formula.

In particular, Meyer takes $S$ large enough that every omitted
prime exceeds the support bound for $f$ and its reciprocal.
Then every summand in $B_{\overline S}(g)$ vanishes.  Formula
(MA-MEY.16) retains the same common-space construction at a
fixed $S$ and exposes the actual prime-carrying remainder.
For $S=\varnothing$ and a test concentrated at $\log2$, the
remainder is exactly $\log2/\sqrt2$, with its reciprocal term
zero.  Thus a prime contribution can be recovered at the
operator level without a bounded narrow-strip M\"obius inverse.

This closes the explicit-commutator and fixed-$S$ common-space
piece.  It does not identify every individual orthogonal
cutoff $M_\rho C_gP_G$ of the two-sided weighted Hilbert model
with an additive $S$-local cutoff kernel.  The equality of
smooth quotient spaces is MT-M24; transporting their
convolution traces to a specified Hilbert quotient requires
the corresponding actual maps and nuclear estimates, under
the precise MT-M26 interior-strip hypothesis.  Likewise,
MA-MEY-002 shows why the compact-cutoff limit cannot simply
be substituted for Meyer's one-sided defect in trace norm.
The source's local additive kernel and the equivariant-Fourier
cutoff in (MA-MEY.15) are actual operators on related local
spaces; equality of their limiting scalar local traces alone
is not an identification of their finite cutoffs.

The geometric interpretation here is the trace of an extension
and its boundary defect.  Meyer explicitly identifies the
cutoff commutator in Lemma~3.2 with the cyclic $1$-cocycle of
a $1$-summable odd Fredholm module in the paragraph following
that lemma, pp.~19--20; the Euler commutator above evaluates
the same one-sided boundary mechanism on the arithmetic
summation operator.  The inherited Grothendieck comparison
with determinants and logarithmic derivatives remains
ALPOGE\_BRIDGES BR-TR-01/02 and F-SM-11/12, cote~1,
scans~079--080 and 111--112.  These archive scans were not
freshly reread in this check.  No identification of this
nonproper analytic limiting operation with a proper algebraic
pushforward in EGA, FGA or SGA is asserted.

\paragraph{Source and check record.}
Meyer, R. (2005). \emph{On a representation of the idele class
group related to primes and zeros of $L$-functions}
(arXiv:math/0311468, version~3, 16 March 2005).
\url{https://arxiv.org/abs/math/0311468}.
The journal version is \emph{Duke Mathematical Journal},
127, 519--595, \url{https://doi.org/10.1215/S0012-7094-04-12734-4};
all page locators here refer to the 53-page arXiv version.
The attached primary \LaTeX\ was used for the complete proofs
of Lemmas~3.2--3.3, the following local principal-value lemmas,
the nuclear-kernel lemma \texttt{lem:mother\_of\_nuclearity},
Proposition/Corollary \texttt{T\_on\_Twist}, and
Lemmas~5.3--5.6 and Theorem~5.8 as used here.  The primary PDF
was consulted for version and page/number locators.  This is
a targeted reading of the complete invoked local/global
proofs, not a claim to have freshly re-proved all upstream
bornological analysis in the paper.  MA-MEY-001--003 have
complete proofs above and individual post-writing check
events in the run record.
